% SGA 3, Expose III, Loop 1 English reconstruction.
% Scope: Corollary 4.31 through Proposition 4.35.1, stopping immediately
% before Proposition 4.36.
% Authority: Polo--Gille Exp3-14oct24.pdf, component-local pp. 71--73
% (combined-reader pp. 181--183; re-edition pp. 171--173; original
% side-pages 153--155). OCR is a locator only. English wording was compared
% with Jacob C. Reinhold's CC BY 4.0 Markdown lineage at commit
% e7a259f3f8608ad3edf9bf6eead3fd504dd2d23e. The Polo--Gille PDF controls
% prose, mathematics, numbering, editorial notes, and the source-derived
% Loop-1 diagram image included below.

\setcounter{footnote}{126}

% Source page: Exp3 local p. 71; combined p. 181; re-edition p. 171;
% original side-page 153.
\medskip
\noindent\SGAInlineRefTarget{sga3:iii:corollary:4-31}\textbf{推论 4.31.}
\textit{设 \(Y\) 如 \SGARefLink{sga3:iii:lemma:4-27}{4.27} 所述.
用 \(X\) 在 \(S\) 上诱导 \(X_{\mathscr J}\) 的单位截面的截面对 \(Y\) 作共轭,
所得共轭子群的集合与下列集合同构:}
\[
 E=\partial C^0(Y_0,L_0^X/L_0^Y)
 =C^0(Y_0,L_0^X/L_0^Y)\big/H^0(Y_0,L_0^X/L_0^Y).
\]

\medskip
\noindent\SGAInlineRefTarget{sga3:iii:corollary:4-32}\textbf{推论 4.32.}
\textit{再设 \(S_0\) 仿射, 且 \(\omega^1_{Y_0/S_0}\) 为有限局部自由模.
\footnote{编者注: 已将 \(\operatorname{Lie}(Y_0/S_0)\)
更正为 \(\omega^1_{Y_0/S_0}\).} 若}
\[
 \mathscr F_0=
 \bigl(\operatorname{Lie}(X_0/S_0)/
       \operatorname{Lie}(Y_0/S_0)\bigr)
 \otimes_{\mathscr O_{S_0}}\mathscr J,
\]
\textit{则 \(E=\Gamma(S_0,\mathscr F_0)/H^0(Y_0,\mathscr F_0)\).}

事实上,\footnote{编者注: 下文对原文作了扩充.}
由于 \(\omega^1_{Y_0/S_0}\) 为有限局部自由模, 而 \(\omega^1_{X_0/S_0}\)
也是如此 (因为已假设 \(X\) 在 \(S\) 上平滑),
\SGARefLink{sga3:iii:remark:0-6}{0.6} 给出
\[
 L_0^Y=\mathbf W\bigl(
   \operatorname{Lie}(Y_0/S_0)\otimes_{\mathscr O_{S_0}}\mathscr J
 \bigr),
 \qquad
 L_0^X=\mathbf W\bigl(
   \operatorname{Lie}(X_0/S_0)\otimes_{\mathscr O_{S_0}}\mathscr J
 \bigr).
\]
另一方面, 由 \SGARefLink{sga3:iii:lemma:4-25}{4.25}, 有
\(Y_0\)-\(\mathscr O_{S_0}\)-模正合序列
\[
 0\longrightarrow \mathscr K\longrightarrow \omega^1_{X_0/S_0}
 \xrightarrow{\varphi}\omega^1_{Y_0/S_0}\longrightarrow0,
 \qquad \mathscr K=\ker(\varphi).
\]
由于 \(\omega^1_{Y_0/S_0}\) 与 \(\omega^1_{X_0/S_0}\) 均为有限局部自由模,
存在局部分裂的正合序列
\[
 0\longrightarrow
 \operatorname{Lie}(Y_0/S_0)\otimes_{\mathscr O_{S_0}}\mathscr J
 \longrightarrow
 \operatorname{Lie}(X_0/S_0)\otimes_{\mathscr O_{S_0}}\mathscr J
 \longrightarrow \mathscr F_0\longrightarrow0.
\]
因此存在 \(Y_0\)-\(\mathscr O_{S_0}\)-模正合序列
\[
 0\longrightarrow L_0^Y\longrightarrow L_0^X
 \longrightarrow\mathbf W(\mathscr F_0).
\]
由证明 \SGARefLink{sga3:iii:corollary:4-31}{4.31} 时所用的论证,
可以将 \(E\) 计算为复合映射
\[
 C^0(Y_0,L_0^X)\xrightarrow{\pi}
 C^0(Y_0,\mathbf W(\mathscr F_0))\xrightarrow{\partial}
 C^1(Y_0,\mathbf W(\mathscr F_0)).
\]
的像. 而 \(\pi\) 就是映射
\[
 \Gamma\bigl(S_0,
   \operatorname{Lie}(X_0/S_0)\otimes_{\mathscr O_{S_0}}\mathscr J
 \bigr)
 \longrightarrow\Gamma(S_0,\mathscr F_0).
\]
由于 \(S_0\) 仿射, \(\pi\) 为满射, 这就证明了所述断言.

\medskip
\noindent\SGAInlineRefTarget{sga3:iii:corollary:4-33}\textbf{推论 4.33.}
\textit{设 \(S,\mathscr I,\mathscr J\) 和 \(X\) 如
\SGARefLink{sga3:iii:theorem:4-21}{4.21} 所述,
并设 \(Y\) 是 \(X\) 的可对角化子群.
假设 \(\omega^1_{Y_0/S_0}\) 为有限局部自由模, 且 \(S_0\) 仿射.
\footnote{编者注: 此处增添了关于 \(\omega^1_{Y_0/S_0}\) 的假设,
并将“\(S\) 仿射”的假设替换为“\(S_0\) 仿射”.}
\(X\) 中那些由 \(X\) 在 \(S\) 上诱导 \(X_{\mathscr J}\) 的单位截面的截面
与 \(Y\) 共轭的子群所成的集合, 与下列集合同构:}
\[
 E=\Gamma\!\left(
 S_0,
 \frac{\operatorname{Lie}(X_0/S_0)}
      {\operatorname{Lie}(X_0/S_0)^{\operatorname{ad}(Y_0)}}
 \otimes_{\Gamma(S_0,\mathscr O_{S_0})}\Gamma(S_0,\mathscr J)
 \right),
\]
\textit{也就是与 \(L_0^X(Y_0)/H^0(Y_0,L_0^X)\) 同构.}
\footnote{编者注: 下文为增补内容; 参见 \SGARefLink{sga3:iii:corollary:4-34}{4.34}.}

事实上, 由 I, \SGARefLink{sga3:i:proposition:4-7-3}{4.7.3}
(参见 \SGARefLink{sga3:iii:corollary:2-13}{2.13}), 写成
\[
 \operatorname{Lie}(X_0/S_0)
 =\operatorname{Lie}(X_0/S_0)^{\operatorname{ad}(Y_0)}
  \oplus\mathscr R.
\]
由于 \(Y_0\) 交换,
\[
 \operatorname{Lie}(Y_0/S_0)
 \subset\operatorname{Lie}(X_0/S_0)^{\operatorname{ad}(Y_0)},
\]
因而
\[
 \begin{aligned}
 \mathscr F_0={}&
 \left(
  \frac{\operatorname{Lie}(X_0/S_0)^{\operatorname{ad}(Y_0)}}
       {\operatorname{Lie}(Y_0/S_0)}
 \right)\otimes_{\mathscr O_{S_0}}\mathscr J
 \ \oplus\ \mathscr R\otimes_{\mathscr O_{S_0}}\mathscr J,\\
 \mathscr F_0^{\operatorname{ad}(Y_0)}={}&
 \left(
  \frac{\operatorname{Lie}(X_0/S_0)^{\operatorname{ad}(Y_0)}}
       {\operatorname{Lie}(Y_0/S_0)}
 \right)\otimes_{\mathscr O_{S_0}}\mathscr J.
 \end{aligned}
\]
因此, 由 \SGARefLink{sga3:iii:corollary:4-32}{4.32}, 有
\(E\simeq
\Gamma(S_0,\mathscr R\otimes_{\mathscr O_{S_0}}\mathscr J)\).
回到 \(\mathscr R\) 的定义即完成证明.

% Source page: Exp3 local p. 72; combined p. 182; re-edition p. 172;
% original side-page 154.
\medskip
\noindent\SGAInlineRefTarget{sga3:iii:corollary:4-34}\textbf{推论 4.34.}
\textit{设 \(S,\mathscr I,\mathscr J\) 和 \(X\) 如
\SGARefLink{sga3:iii:theorem:4-21}{4.21} 所述,
并设 \(Y\) 是 \(X\) 的可对角化子群.
假设 \(\omega^1_{Y_0/S_0}\) 为有限局部自由模, 且 \(S_0\) 仿射.
\footnote{编者注: 参见编者注 129.}
若 \(x\in X(S)\) 诱导 \(X_{\mathscr J}\) 的单位截面并正规化 \(Y\),
则它中心化 \(Y\).}

比较前一推论与 \SGARefLink{sga3:iii:corollary:2-14}{2.14}, 立即得到这一结论.
事实上, \SGARefLink{sga3:iii:corollary:4-33}{4.33} 表明,
\(C^0(Y_0,L_0^X)\) 中整体保持 \(Y\) 的元素正是
\(H^0(Y_0,L_0^X)\) 的元素; 而 \SGARefLink{sga3:iii:corollary:2-14}{2.14} 表明,
这些元素恰好是平凡地作用于典范浸入 \(Y\longrightarrow X\) 的元素.

\medskip
\noindent\SGAInlineRefTarget{sga3:iii:paragraph:4-35}\textbf{4.35. ---}
回到 \SGARefLink{sga3:iii:theorem:4-21}{4.21} 的一般情形,
并假设 \(Y_{\mathscr J}\) 在 \(S_{\mathscr J}\) 上平滑.
由 \SGARefLink{sga3:iii:lemma:4-25:item:ii}{4.25(ii)}, 此时有
\(Y_0\)-\(\mathscr O_{S_0}\)-模正合序列
\[
 \tag{*}\label{sga3:iii:paragraph:4-35:equation:star}
 0\longrightarrow\operatorname{Lie}(Y_0/S_0)
 \longrightarrow\operatorname{Lie}(X_0/S_0)
 \longrightarrow\mathbf n_{Y_0/X_0}^{\vee}\longrightarrow0,
\]
且其中各模均为有限局部自由 \(\mathscr O_{S_0}\)-模.

另一方面, 由 SGA~1, II, \SGARefLink{sga3:ii:definition:4-10}{4.10},
\(X\) 中每个提升 \(Y_{\mathscr J}\) 且在 \(S\) 上平坦的子概形 \(Y\),
都在 \(S\) 上平滑.\footnote{编者注: 下文及
\SGARefLink{sga3:iii:proposition:4-35-1}{命题 4.35.1} 为增补内容;
原文中已隐含这些内容. 参见 4.38(5).}
再设 \(S_0\) 和 \(Y_{\mathscr J}\) 仿射.
由于 \(\mathbf n_{Y_0/X_0}^{\vee}\) 为局部自由 \(\mathscr O_{S_0}\)-模,
序列 \((*)\) 先在 \(\mathscr O_{S_0}\) 上与 \(\mathscr J\) 作张量积,
再拉回到 \(Y_0^n\) 后仍保持正合. 由于各 \(Y_0^n\) 均仿射,
这给出一个交换群复形的正合序列
\[
 0\longrightarrow C^*(Y_0,L_0^Y)
 \longrightarrow C^*(Y_0,L_0^X)
 \xrightarrow{d^*}C^*(Y_0,N_0)\longrightarrow0,
\]
特别地, 得到一个各行正合的交换图:

% Raster stand-in for Loop 1; replace by a native TeX diagram in Loop 2.
% Authority: Exp3 local p. 72; combined reader p. 182; re-edition p. 172;
% original side-page 154.
\begin{center}
 \SGAThreeNativeInput[0.94\textwidth]
   {figures/exp3/diagram_062_exp3_p072_435_cochain_diagram.tex}
\end{center}

现设 \(Y\) 和 \(Y'\) 是 \(X\) 中两个提升 \(Y_{\mathscr J}\) 且在 \(S\)
上平坦, 因而平滑的子群. 由于 \(Y_{\mathscr J}\) 仿射,
\SGARefLink{sga3:iii:proposition:0-15}{0.15} 表明,
\(Y\) 和 \(Y'\) 作为提升 \(Y_{\mathscr J}\) 的概形同构;
也就是说, 存在一个 \(S\)-概形同构
\[
 f:Y\xrightarrow{\sim}Y'
\]

% Source page: Exp3 local p. 73; combined p. 183; re-edition p. 173;
% original side-page 155.
在 \(Y_{\mathscr J}\) 上诱导恒等映射.
一方面, 由 \SGARefLink{sga3:iii:proposition:1-2-4}{1.2.4},
\(f\) 定义了一个元素 \(a\in C^1(Y_0,L_0^X)\), 使得对每个
\(y\in Y(S')\), \(S'\to S\), 有
\[
 f(y)=a(y_0)y
\]
且由 \SGARefLink{sga3:iii:lemma:4-27:item:i}{4.27(i)}, 有
\[
 d^1(a)=d(Y,Y').
\]
此外, 由于 \(Y\) 和 \(Y'\) 是 \(X\) 的子群,
最后这个元素属于 \(Z^1(Y_0,N_0)\) (参见 \SGARefLink{sga3:iii:theorem:4-21}{4.21}).
因此 \(\partial a\) 是 \(Z^2(Y_0,L_0^Y)\) 的元素,
其在 \(H^2(Y_0,L_0^Y)\) 中的像 \(\overline{\partial a}\)
仅依赖于类 \(\overline{d(Y,Y')}\in H^1(Y_0,N_0)\).
这正是连接映射
\[
 \partial^1:H^1(Y_0,N_0)\longrightarrow H^2(Y_0,L_0^Y),
\]
的定义, 因而
\[
 \partial^1\bigl(\overline{d(Y,Y')}\bigr)
 =\overline{\partial a}.
\]

另一方面, 通过 \(f\) 转移 \(Y'\) 的群结构, 以 \(Y_1\) 表示所得的群,
所以其底概形为 \(Y\). 其群律 \(\mu_1\) 对每个 \(S'\to S\) 和
\(x,y\in Y(S')\) 由下式定义:
\[
 \mu_1(x,y)=f^{-1}\bigl(f(x)f(y)\bigr).
\]
由 \SGARefLink{sga3:iii:remark:3-5-1}{3.5.1}, \(Y_1\) 定义了一个余循环
\[
 \delta(Y,Y_1)\in
 Z^2\bigl(Y_0,\operatorname{Lie}(Y_0/S_0)\bigr)
\]
使得对每个 \(S'\to S\) 及 \(x,y\in Y(S')\),
\[
 \delta(Y,Y_1)(x_0,y_0)xy
 =\mu_1(x,y)=f^{-1}\bigl(f(x)f(y)\bigr).
\]
置 \(b=\delta(Y,Y_1)\). 对每个 \(S'\to S\) 和 \(x,y\in Y(S')\),
有 \((b(x_0,y_0)xy)_0=x_0y_0\). 因而, 一方面,
\(f(b(x_0,y_0)xy)\) 等于
\[
 a(x_0y_0)b(x_0,y_0)xy,
\]
另一方面又等于
\[
 f(x)f(y)=a(x_0)x\,a(y_0)y
 =a(x_0)\operatorname{Ad}(x_0)\bigl(a(y_0)\bigr)xy.
\]
比较这两个表达式可知, \(b(x_0,y_0)\) 等于
\[
 \begin{aligned}
 &a(x_0y_0)^{-1}a(x_0)
   \operatorname{Ad}(x_0)\bigl(a(y_0)\bigr)\\
 &\qquad={}
 \operatorname{Ad}(x_0)\bigl(a(y_0)\bigr)-a(x_0y_0)+a(x_0)
 =(\partial a)(x_0,y_0).
 \end{aligned}
\]
因此 \(\delta(Y,Y_1)=\partial a\), 从而我们证明了:

\medskip
\noindent\SGAInlineRefTarget{sga3:iii:proposition:4-35-1}\textbf{命题 4.35.1.}\footnotemark[132]
\textit{在 \SGARefLink{sga3:iii:theorem:4-21}{4.21} 的假设下,
再设 \(S_0\) 仿射, 且 \(Y_{\mathscr J}\) 在 \(S_{\mathscr J}\) 上平滑并仿射.
设 \(Y,Y'\) 是 \(X\) 的两个子群, 它们提升 \(Y_{\mathscr J}\) 且在 \(S\)
上平坦 (因而平滑). 设 \(f:Y\xrightarrow{\sim}Y'\) 是在
\(Y_{\mathscr J}\) 上诱导恒等映射的 \(S\)-概形同构,
并以 \(Y_1\) 表示通过 \(f\) 转移 \(Y'\) 的群结构而得到的群. 则}
\[
 \partial^1\bigl(\overline{d(Y,Y')}\bigr)
 =\overline{\delta(Y,Y_1)}.
\]

% Continue with Proposition 4.36 on Exp3 local p. 73 / combined p. 183 /
% re-edition p. 173 / original side-page 155.
